% ==============================================================================
% CHƯƠNG 4: KẾT LUẬN VÀ HƯỚNG PHÁT TRIỂN
% ==============================================================================
\chapter{KẾT LUẬN VÀ HƯỚNG PHÁT TRIỂN}
\label{chap:ket_luan}

\section{Đánh giá những kết quả đạt được}
\label{sec:ket_qua_dat_duoc}

Sau thời gian nghiêm túc nghiên cứu lý thuyết, khảo sát nghiệp vụ thực tế và triển khai lập trình, đề tài \textbf{''Xây dựng website POS Quản lý bán hàng trên nền tảng Spring Boot và ReactJS''} đã hoàn thành xuất sắc toàn bộ các mục tiêu đặt ra ban đầu, đáp ứng đầy đủ các yêu cầu khắt khe của môn học \textit{Phát triển phần mềm mã nguồn mở} tại Trường Đại học Giao thông Vận tải.

Cụ thể, các kết quả chính mà nhóm đã đạt được bao gồm:

\subsection{Về mặt nghiên cứu lý thuyết và phương pháp luận}
\begin{itemize}[leftmargin=1.5cm]
    \item \textbf{Làm chủ kiến trúc phân tán hiện đại:} Nắm vững và vận dụng thành thạo mô hình kiến trúc phân tầng (Multi-tier Architecture) tách biệt giữa giao diện máy khách (React 19 SPA) và máy chủ dịch vụ (Spring Boot 4.x RESTful API), kết hợp hạ tầng container hóa Docker và Nginx Reverse Proxy.
    \item \textbf{Chuẩn hóa phân tích thiết kế hệ thống:} Xây dựng hệ thống biểu đồ phân tích bài bản từ Biểu đồ phân rã chức năng (BFD), mô hình Use Case tổng quát đến 6 bảng đặc tả ca sử dụng chi tiết chuẩn học thuật; mô hình hóa quy trình vận hành qua các biểu đồ Hoạt động, Tuần tự, Biểu đồ lớp và Biểu đồ triển khai phần cứng.
    \item \textbf{Thiết kế cơ sở dữ liệu quan hệ chuẩn mực:} Xây dựng mô hình quan hệ thực thể (ERD) đạt chuẩn dạng chuẩn 3 (3NF) với 13 bảng dữ liệu thực tế, giải quyết triệt để quan hệ nhiều - nhiều giữa sản phẩm và nhóm tùy chọn thông qua bảng trung gian, áp dụng công cụ Flyway Migration quản lý phiên bản schema nghiêm ngặt thay vì cấu hình tự sinh bảng rủi ro.
    \item \textbf{Thiết lập các chuẩn an ninh và kiến trúc cốt lõi (ADR):}
    \begin{itemize}
        \item \textit{Bảo mật hai lớp token (ADR-0007):} Phối hợp Access Token trong RAM và Refresh Token trong Cookie HttpOnly với cơ chế Token Rotation và Token Family ngăn chặn hoàn toàn nguy cơ rò rỉ phiên và tấn công Replay Attack.
        \item \textit{Quản lý tồn kho theo sổ cái kế toán (ADR-0004):} Mọi biến động kho đều được ghi nhận bằng các giao dịch độc lập (\texttt{IN}, \texttt{OUT}, \texttt{ADJUSTMENT}), đảm bảo tính minh bạch và truy vết lịch sử kho hàng.
        \item \textit{Bảo vệ toàn vẹn tài chính (ADR-0006):} Cấm xóa đơn hàng đã hoàn tất, tự động kích hoạt quy trình giao dịch bù trừ toàn diện (hoàn kho, hoàn lượt voucher, hoàn CRM) khi hủy đơn chờ.
        \item \textit{Khuyến mãi không xếp chồng (ADR-0003):} Đánh giá tự động và áp dụng duy nhất 1 chương trình ưu đãi tối ưu nhất cho khách hàng.
    \end{itemize}
\end{itemize}

\subsection{Về mặt thực nghiệm và sản phẩm ứng dụng}
\begin{itemize}[leftmargin=1.5cm]
    \item \textbf{Hệ thống giao diện POS trực quan, mượt mà:} Xây dựng hoàn chỉnh 11 màn hình giao diện người dùng trên nền tảng React 19 và Tailwind CSS v4, tối ưu hóa cho màn hình cảm ứng tại quầy; thời gian phản hồi thao tác thêm món và cập nhật giỏ hàng dưới $50$ mili-giây.
    \item \textbf{Tích hợp thanh toán trực tuyến VietQR động PayOS:} Kết nối thành công cổng PayOS qua bộ SDK Java chính thức, sinh mã VietQR động theo chuẩn NAPAS có khóa số tiền và nội dung đơn; xử lý Webhook tự động cập nhật trạng thái đơn hàng kèm xác thực chữ ký điện tử HMAC-SHA256 Checksum; hỗ trợ quy trình chuyển đổi sang tiền mặt linh hoạt khi gặp sự cố mạng.
    \item \textbf{Tích hợp lưu trữ đám mây AWS S3:} Lưu trữ hình ảnh sản phẩm và danh mục phân tán trên hạ tầng Amazon S3 thông qua AWS SDK v2, tự động dọn dẹp tệp ảnh cũ khi cập nhật dữ liệu.
    \item \textbf{Mẫu hóa đơn in nhiệt 80mm chuẩn hóa:} Hỗ trợ in trực tiếp từ trình duyệt qua lệnh \texttt{window.print()} với định dạng CSS \texttt{@media print} sắc nét.
    \item \textbf{Chất lượng kiểm thử tuyệt đối:} Vượt qua 100\% các ca kiểm thử chức năng, toàn vẹn kế toán và an toàn bảo mật trong ma trận 15 kịch bản kiểm thử (15/15 ca PASS); thời gian đáp ứng trung bình của máy chủ API đạt từ $120$ms đến $250$ms.
\end{itemize}

% ==============================================================================
\section{Những hạn chế còn tồn tại}
\label{sec:han_che}

Mặc dù hệ thống đã đáp ứng đầy đủ các yêu cầu cốt lõi và vận hành ổn định trong môi trường thực nghiệm, song do giới hạn về mặt thời gian và phạm vi của một đề tài bài tập lớn, ứng dụng vẫn còn một số điểm hạn chế nhất định:
\begin{enumerate}[label=\arabic*., leftmargin=1.5cm]
    \item \textbf{Chưa hỗ trợ chế độ ngoại tuyến hoàn chỉnh (Offline-First):} Hệ thống hiện tại vận hành theo cơ sở dữ liệu tập trung trên máy chủ. Trong trường hợp điểm bán lẻ bị mất hoàn toàn kết nối Internet, thu ngân chỉ có thể duy trì việc bán hàng tạm thời nếu đã nạp sẵn dữ liệu trang trước đó, nhưng việc lưu đơn và cập nhật tồn kho lên cơ sở dữ liệu trung tâm vẫn đòi hỏi kết nối mạng khả dụng.
    \item \textbf{Phụ thuộc vào tính sẵn sàng của dịch vụ bên thứ ba:} Luồng thanh toán chuyển khoản ngân hàng phụ thuộc vào thời gian hoạt động (Uptime) của cổng PayOS và hệ thống Napas/Ngân hàng thụ hưởng; luồng hiển thị hình ảnh phụ thuộc vào kết nối tới cụm máy chủ Amazon S3 khu vực Singapore.
    \item \textbf{Giao tiếp phần cứng ngoại vi còn mang tính mô phỏng:} Việc kết nối với máy quét mã vạch hiện tại được xử lý thông qua cơ chế bàn phím ảo (Keyboard Wedge input) thay vì giao tiếp trực tiếp qua cổng COM/Serial phần cứng cấp thấp; mẫu in hóa đơn phụ thuộc vào trình điều khiển in của trình duyệt thay vì điều khiển trực tiếp bằng tập lệnh ESC/POS.
\end{enumerate}

% ==============================================================================
\section{Hướng phát triển trong tương lai}
\label{sec:huong_phat_trien}

Để phát triển hệ thống POS từ quy mô một đồ án môn học trở thành một giải pháp phần mềm thương mại hoàn chỉnh có khả năng cạnh tranh cao trên thị trường, nhóm định hướng các giai đoạn nâng cấp tiếp theo:
\begin{enumerate}[label=\arabic*., leftmargin=1.5cm]
    \item \textbf{Phát triển ứng dụng di động cho Thủ kho (Mobile App):} Xây dựng ứng dụng di động đa nền tảng bằng React Native dành riêng cho nhân viên kho, hỗ trợ quét mã vạch SKU bằng camera điện thoại để thực hiện kiểm kê, nhập hàng và xuất hủy nhanh chóng ngay tại kệ hàng.
    \item \textbf{Triển khai kiến trúc Progressive Web App (PWA) \& Đồng bộ Offline:} Tích hợp Service Workers và cơ sở dữ liệu cục bộ IndexedDB trên trình duyệt để lưu trữ tạm thời các đơn hàng phát sinh khi mất mạng, sau đó tự động thực hiện cơ chế đồng bộ nền (Background Sync) lên máy chủ khi kết nối Internet được khôi phục.
    \item \textbf{Ứng dụng Trí tuệ nhân tạo (AI) trong quản trị bán lẻ:}
    \begin{itemize}
        \item \textit{Dự báo nhu cầu nhập hàng (Demand Forecasting):} Phân tích chuỗi thời gian lịch sử bán hàng kết hợp các yếu tố mùa vụ, ngày lễ để dự báo chính xác số lượng hàng hóa cần nhập trong tuần/tháng tiếp theo, giảm thiểu chi phí lưu kho.
        \item \textit{Hệ thống gợi ý bán chéo thông minh (Smart Cross-selling):} Phân tích quy luật giỏ hàng (Market Basket Analysis) để gợi ý thu ngân mời khách hàng mua kèm các món phù hợp (ví dụ: gợi ý thêm Topping trân châu khi khách gọi trà sữa).
    \end{itemize}
    \item \textbf{Mở rộng mô hình Chuỗi cửa hàng đa chi nhánh (Multi-branch POS):} Nâng cấp cơ sở dữ liệu hỗ trợ nhiều kho hàng độc lập theo từng chi nhánh, cho phép điều chuyển hàng hóa nội bộ giữa các cửa hàng và phân quyền quản lý theo từng cụm chi nhánh.
    \item \textbf{Kết nối trực tiếp thiết bị ngoại vi qua chuẩn WebUSB / WebSerial:} Gửi trực tiếp các lệnh điều khiển máy in nhiệt ESC/POS (cắt giấy tự động, mở ngăn kéo đựng tiền - Cash Drawer) và đọc dữ liệu cân điện tử từ cổng nối tiếp mà không cần cài đặt driver trung gian.
\end{enumerate}

% ==============================================================================
\section{Lời kết}
Quá trình thực hiện đề tài \textbf{''Xây dựng website POS Quản lý bán hàng trên nền tảng Spring Boot và ReactJS''} không chỉ giúp nhóm sinh viên củng cố vững chắc các kiến thức lý thuyết đã học về công nghệ phần mềm, lập trình hướng đối tượng, phân tích thiết kế hệ thống và cơ sở dữ liệu, mà còn mang lại những kinh nghiệm thực tiễn vô cùng quý báu về việc áp dụng các công nghệ mã nguồn mở tiên tiến nhất vào giải quyết một bài toán nghiệp vụ kinh doanh thực tế. Hệ thống đã chứng minh được tính khả thi cao, sự ổn định trong vận hành và tiềm năng mở rộng to lớn để trở thành một sản phẩm phần mềm thương mại hoàn chỉnh trong tương lai.
